%=====================================================================
%  IEEE journal manuscript. Compile with pdfLaTeX + BibTeX
%  Order:  pdfLaTeX -> BibTeX -> pdfLaTeX -> pdfLaTeX
%  Multimodal polysomnography: joint staging + respiratory-event detection.
%=====================================================================
\documentclass[journal]{IEEEtran}
\IEEEoverridecommandlockouts
\usepackage{cite}
\usepackage{amsmath,amssymb,amsfonts}
\usepackage{algorithmic}
\usepackage{graphicx}
\usepackage{textcomp}
\usepackage{xcolor}
\usepackage{booktabs}
\usepackage{array}
\usepackage{multirow}
\usepackage{url}
\def\BibTeX{{\rm B\kern-.05em{\sc i\kern-.025em b}\kern-.08em T\kern-.1667em\lower.7ex\hbox{E}\kern-.125emX}}

% Headline numbers (single source of truth, results/mm_final_cv.json).
\newcommand{\mmAcc}{0.721}   % multimodal staging acc (+HMM)
\newcommand{\mmMf}{0.645}    % multimodal staging macro-F1 (+HMM)
\newcommand{\mmK}{0.606}     % multimodal staging kappa (+HMM)
\newcommand{\mmAuc}{0.700}   % multimodal apnea AUC
\newcommand{\mmAp}{0.333}    % multimodal apnea average precision
\newcommand{\eegAuc}{0.655}  % EEG-only apnea AUC (cortical arousals)
\newcommand{\noflowAuc}{0.707} % apnea AUC with airflow removed

\begin{document}

\title{Multimodal Polysomnography for Joint Sleep Staging and Respiratory-Event Detection in Subacute Ischemic Stroke}

\author{Md Wahiduzzaman Sua, Esm-e Moula Chowdhury Abha, and Md Imtiaj Alam Sajin%
\thanks{The authors are with the Department of Computer Science and Engineering,
American International University-Bangladesh, Dhaka, Bangladesh
(e-mail: 26-94088-2@student.aiub.edu; 26-94089-2@student.aiub.edu;
26-94090-2@student.aiub.edu).}}

\markboth{Manuscript, July 2026}%
{Sua \MakeLowercase{\textit{et al.}}: Multimodal Polysomnography in Ischemic Stroke}

\maketitle

%=====================================================================
\begin{abstract}
A clinical polysomnogram records far more than the brain. Alongside the
electroencephalogram it captures eye movements, muscle tone, the electrocardiogram,
nasal and oral airflow, thoracic and abdominal breathing effort, and blood oxygen
saturation. In patients recovering from ischemic stroke, a group in which breathing is
disturbed during sleep at high rates, these cardiorespiratory signals carry the disease
itself. We build one model that reads the full polysomnogram and produces two clinical
outputs at once: the sleep stage of every 30-second epoch and a per-epoch label for
respiratory events. The model encodes the neural and the cardiorespiratory signals in two
parallel streams, lets the two streams exchange information, follows the night with a
bidirectional recurrent decoder, and drives a staging head and a respiratory head. A
direct connection carries the cardiorespiratory signal to the respiratory head so that
oxygen desaturation reaches the decision at full strength. On iSLEEPS, the first public
polysomnography corpus of subacute ischemic stroke (99 patients), evaluated with
patient-independent ten-fold cross-validation, the single model stages sleep at $\mmAcc$
accuracy with Cohen's $\kappa$ of $\mmK$ and detects respiratory events at $\mmAuc$ area
under the ROC curve. The respiratory read-out is grounded in real pathophysiology: when we
remove the airflow channel from which the events were originally scored, detection holds
at $\noflowAuc$, because the model reads oxygen desaturation, breathing effort, cardiac
variability, and the cortical arousals that follow an obstructed breath. Each modality
contributes where its physiology lives. The cortical channels carry sleep staging, and the
cardiorespiratory channels give the model a second clinical capability that an
electroencephalogram alone cannot provide. Code and features are released.
\end{abstract}

\begin{IEEEkeywords}
Polysomnography, multimodal learning, sleep staging, sleep apnea, ischemic stroke,
cardiorespiratory signals, multi-task learning.
\end{IEEEkeywords}

%=====================================================================
\section{Introduction}
\IEEEPARstart{S}{leep} is scored in 30-second epochs into five stages, Wake (W), N1, N2,
N3, and rapid-eye-movement (REM), following the American Academy of Sleep Medicine (AASM)
rules~\cite{aasm2023manual}. The overwhelming majority of automated methods read a single
electroencephalogram (EEG) channel and ask one question: how accurately can the five stages
be classified? Two decades of work on that question have produced convolutional, recurrent,
attention, and self-supervised EEG models~\cite{supratak2017deepsleepnet,phan2019seqsleepnet,
eldele2021attnsleep,phan2021xsleepnet,labram2024,cbramod2025,eegpt2024}.

A clinical polysomnogram (PSG), however, is a recording of the whole sleeping body. It
carries the EEG, eye movements, muscle tone, the electrocardiogram (ECG), nasal and oral
airflow, thoracic and abdominal breathing effort, and peripheral oxygen saturation
(SpO$_2$). These cardiorespiratory channels are where sleep-disordered breathing writes
itself: a paused breath, a chest that keeps straining against a closed airway, a fall in
blood oxygen, and the arousal that rescues the breath. A model that reads only the EEG
cannot see any of this, and so cannot report on the disorder that most damages a patient's
sleep. We use the full polysomnogram, and we ask a question that the staging race does not:
what does a model gain, and on which clinical task, when it reads the whole recording rather
than the brain alone.

We study this on iSLEEPS~\cite{maiti2026isleeps}, released in 2026 as the first public PSG
corpus dedicated to subacute ischemic stroke. The cohort is well matched to the question.
Sleep-disordered breathing is common in the weeks after a stroke and is tied to worse
recovery~\cite{bassetti2006sleep,denis2024sleep}, and in this corpus the median patient
carries a scored respiratory event in $13\%$ of epochs, with $77$ of $96$ patients above
$5\%$ (Fig.~\ref{fig:prev}). The cardiorespiratory signal is therefore abundant and
clinically meaningful in exactly the population where breathing matters most.

We present a two-stream, multi-task model that treats staging and respiratory-event
detection as one joint problem. A neural stream and a cardiorespiratory stream are encoded
in parallel and combined, a bidirectional recurrent decoder reads the resulting sequence
across the night, and two heads emit the sleep stage and a per-epoch respiratory label. The
model is compact by design, which suits a cohort of 99 patients, and every part is built to
be measured rather than assumed. Our contributions are:

\begin{itemize}
\item \textbf{One model for two clinical read-outs from the full polysomnogram.} It stages
sleep at $\mmAcc$ accuracy ($\kappa$ $\mmK$) and, from the same forward pass, detects
respiratory events at $\mmAuc$ AUC, with a direct cardiorespiratory path to the respiratory
head so oxygen-desaturation cues reach the decision undiminished (Section~\ref{sec:method}).
\item \textbf{A clear reading of where each modality contributes.} Under a
patient-independent ten-fold protocol, the cardiorespiratory channels lift respiratory-event
detection from $\eegAuc$, the level reachable from cortical arousals alone, to $\mmAuc$
AUC, while the neural channels carry sleep staging. Each signal earns its place on the task
its physiology governs (Section~\ref{sec:contrib}).
\item \textbf{A respiratory read-out grounded in physiology, not in the labeling signal.}
When the airflow channel from which the events were scored is removed, detection holds at
$\noflowAuc$ AUC, because the decision rests on oxygen desaturation, breathing effort,
cardiac variability, and cortical arousals (Section~\ref{sec:noflow}).
\end{itemize}

%=====================================================================
\section{Related Work}
\label{sec:related}

\subsection{From single-channel to multi-channel sleep staging}
Automated stagers began with one EEG channel~\cite{supratak2017deepsleepnet,
phan2019seqsleepnet,eldele2021attnsleep} and now read several PSG channels.
U-Sleep~\cite{perslev2021usleep} stages from arbitrary EEG and EOG derivations across many
cohorts, and multi-channel models add EOG, EMG, and in some cases
ECG~\cite{phan2021xsleepnet,biot2023,sleepfm2026}. The eye and muscle channels help most
with REM and Wake, which a single EEG channel resolves poorly. Our work reads the
cardiorespiratory channels as well, and reports both what they give and where they give it.

\subsection{Detection of sleep-disordered breathing}
A separate line of work detects apnea and hypopnea from breathing signals. Airflow and
effort shape, oximetry, and ECG-derived heart-rate variability are all informative, and
oxygen desaturation is a defining marker of an event~\cite{jing2025cpap}. These detectors
usually run at fine time resolution on the raw breathing waveforms and are trained on that
task alone. We detect respiratory events on the staging epoch grid, jointly with staging,
from a shared representation, so that one model produces both clinical outputs and the two
tasks support each other during training.

\subsection{Sleep and breathing after stroke}
Sleep is both disturbed by stroke and predictive of recovery, and sleep-disordered breathing
is common in the subacute phase and linked to poorer functional
outcome~\cite{bassetti2006sleep,denis2024sleep}. Focal injury also changes the micro-events
that scoring relies on, the spindles, slow waves, and muscle atonia~\cite{gottselig2002stroke,
fernandez2019spindles,simpson2023laterality}. A model that reports on breathing as well as
stage is therefore of direct clinical value in this population.

%=====================================================================
\section{The polysomnogram we model}
\label{sec:data}
iSLEEPS~\cite{maiti2026isleeps} contains overnight PSG from 100 subacute ischemic-stroke
patients (mean age $50.5$; 77 male, 23 female), scored by experts into the five AASM stages.
We exclude one recording after an audit: the data records of SN28 are byte-identical to
those of SN15, a duplication consistent with the anonymization and re-numbering the dataset
authors describe, and keeping it would place a training subject into the test fold in nine of
ten fold assignments. We therefore use $N=99$ ($96$ with complete cardiorespiratory
channels), a total of $89{,}532$ epochs. The stage distribution (Table~\ref{tab:dist}) is
strongly imbalanced, with N2 dominant and N1 and N3 scarce, which is the clinical norm.

\begin{table}[t]
\caption{Stage distribution and respiratory-event prevalence in the corpus we model ($N=99$).}
\label{tab:dist}
\centering
\begin{tabular}{lccccc}
\toprule
Stage & W & N1 & N2 & N3 & R \\
\midrule
Share (\%) & 26.6 & 10.2 & 42.3 & 8.9 & 12.1 \\
\bottomrule
\end{tabular}
\\[2pt]
\footnotesize Respiratory events: $16.0\%$ of epochs positive; median $13\%$ per patient
(range $0$ to $60.5\%$); $77$ of $96$ patients above $5\%$.
\end{table}

\begin{figure}[t]
\centering
\includegraphics[width=0.86\columnwidth]{figures/fig_mm_apnea_prev.pdf}
\caption{Sleep-disordered-breathing burden across the cohort. Each bar counts patients by
the fraction of epochs that carry a scored respiratory event. The median patient sits at
$13\%$ and many exceed $30\%$, so the cardiorespiratory signal is abundant.}
\label{fig:prev}
\end{figure}

\paragraph{Channels we read, and channels we do not}
A raw recording holds 22 traces. We read the 14 that carry sleep-stage or breathing
information. The \emph{neural stream} takes the four EEG derivations (C4:M1, C3:M2, O2:M1,
O1:M2), the two EOG channels (E1:M2, E2:M2), and the submental EMG. The
\emph{cardiorespiratory stream} takes the ECG, nasal and oral airflow, thoracic effort,
abdominal effort, the summed effort trace, SpO$_2$, and pulse. We leave out eight traces for
stated reasons: periodic limb movement records a different disorder; snore is scored from the
same airway sound as the respiratory events and would make the second task circular;
plethysmography is already summarized by the SpO$_2$ and pulse channels we keep; body
position is coarse and slow; and the ambient-light and battery traces are device
housekeeping, not physiology. All neural channels are resampled to $100$\,Hz and the
cardiorespiratory channels to $25$\,Hz, and both are read in physical units, which matters
because SpO$_2$ near $90\%$, pulse near $78$\,bpm, and microvolt EEG live on very different
scales.

\paragraph{Respiratory-event labels}
The second task is supervised by the corpus's scored respiratory events, mapped onto the
$30$\,s epoch grid: an epoch is positive when it overlaps a scored event.
Section~\ref{sec:noflow} shows that our model does not simply re-read the airflow trace these
labels came from.

%=====================================================================
\section{Model}
\label{sec:method}
We map each night to two per-epoch feature sequences and learn a joint model
$f_\theta:(\mathbf{f}^{\mathrm{eeg}}_{1:T},\mathbf{f}^{\mathrm{car}}_{1:T})\rightarrow
(y^{\mathrm{stg}}_{1:T},y^{\mathrm{apn}}_{1:T})$ that generalizes to unseen patients
(Fig.~\ref{fig:arch}).

\begin{figure*}[t]
\centering
\includegraphics[width=\textwidth]{figures/fig_architecture.pdf}
\caption{The two-stream, multi-task model. (A) The neural features and the cardiorespiratory
features are encoded in parallel, combined by a cross-modal fusion block, and read across the
night by a bidirectional LSTM. A staging head reads the recurrent state; a respiratory head
reads the recurrent state together with a direct copy of the cardiorespiratory embedding, so
oxygen-desaturation and effort cues reach the respiratory decision without passing through the
staging-oriented fusion. (B) Detail of the feature encoder and of the fusion block, which
tokenizes the two modalities, runs four-head attention, and combines the result with a
residual connection.}
\label{fig:arch}
\end{figure*}

\subsection{Two physiological feature sets}
Each epoch's neural channels are summarized by $188$ features that reflect how a technician
reads sleep: absolute and relative band powers in delta, theta, alpha, sigma, and beta from
the Welch spectrum, spectral entropy and edge frequency, the three Hjorth
descriptors~\cite{hjorth1970eeg}, time-domain statistics, spindle density and envelope
moments from the sigma-band analytic signal, slow-wave amplitudes, ocular-movement energy on
the EOG, and tonic EMG level. From the seven cardiorespiratory channels we compute $14$
features chosen to describe the mechanics of a respiratory event: SpO$_2$ mean, minimum,
variability, and desaturation depth; pulse rate and its variability; ECG amplitude and
line-length; airflow amplitude and line-length; thoracic, abdominal, and summed effort
amplitudes; and thoraco-abdominal asynchrony, the loss of coordination between chest and
abdomen that marks an obstructed breath. Both sets are standardized per patient, which
removes person-to-person offsets so the model learns across people.

\subsection{Two-stream encoding and cross-modal fusion}
Each stream is a two-layer perceptron with LayerNorm that maps the neural features to
$e_t\in\mathbb{R}^{128}$ and the cardiorespiratory features to $c_t\in\mathbb{R}^{64}$. The
fusion block treats $e_t$ and $c_t$ as two tokens, adds learned modality-type embeddings,
runs four-head attention so each modality can read the other, and combines the two refined
tokens with a residual connection into $z_t\in\mathbb{R}^{128}$. We also test a plain
concatenation in place of attention; the two perform alike, and we report both.

\subsection{Temporal decoder and two heads}
Sleep is sequential, and respiratory events cluster, so the fused vectors of a $20$-epoch
window ($10$ minutes) are read by a two-layer bidirectional LSTM with $128$ hidden units per
direction, giving $h_t\in\mathbb{R}^{256}$. The staging head is a linear map to five classes.
The respiratory head reads $h_t$ together with a direct copy of the per-epoch
cardiorespiratory embedding $c_t$. This direct connection is deliberate: without it, the
respiratory decision must survive a fusion trained mostly for staging, and the fall in oxygen
that defines an event is weakened along the way. The full model has $0.86$ million parameters.

\subsection{Training and decoding}
The two tasks are trained together with equal weight,
\begin{equation}
\mathcal{L}=\mathcal{L}^{\mathrm{stg}}_{\mathrm{CE}}+\lambda\,\mathcal{L}^{\mathrm{apn}}_{\mathrm{BCE}},
\qquad \lambda=1 .
\end{equation}
Two choices carry the staging accuracy. The staging cross-entropy is weighted by the square
root of inverse class frequency, which corrects the scarce N1 and N3 stages without giving up
overall accuracy, and at inference the staging posteriors are smoothed by a hidden-Markov
Viterbi pass whose transitions are counted from the training labels~\cite{chen2024hmm}, which
restores the natural persistence of sleep. The respiratory head is positive-weighted to match
the $16\%$ event rate. Optimization uses AdamW with cosine annealing and early stopping on a
held-out set of patients.

%=====================================================================
\section{Results}
\label{sec:results}
\subsection{Protocol}
All results use ten-fold, patient-independent cross-validation, so no patient appears in both
training and test. Within each training split we hold out a disjoint set of patients for early
stopping. We report staging accuracy, macro-F1, and Cohen's $\kappa$, and for the respiratory
task the area under the ROC curve (AUC) and average precision (AP), which suits the $16\%$
event rate.

\subsection{Two clinical read-outs from one model}
The model stages sleep at $\mmAcc$ accuracy, $\mmMf$ macro-F1, and $\mmK$ $\kappa$, and detects
respiratory events at $\mmAuc$ AUC with $\mmAp$ average precision (Table~\ref{tab:main}). One
compact model delivers both outputs on unseen patients from a single pass over the recording.
Staging at this level on a stroke cohort is a strong result: the injured brain distorts the
very spindles, slow waves, and atonia that scoring depends on, and every epoch is judged on a
patient the model has never seen.

\begin{table}[t]
\caption{The joint multimodal model, ten-fold patient-independent. Staging is reported with
Viterbi smoothing; the respiratory task by AUC and average precision. Fold standard deviations
in parentheses.}
\label{tab:main}
\centering
\setlength{\tabcolsep}{4pt}
\begin{tabular}{lccc}
\toprule
Output & \multicolumn{3}{c}{Metric}\\
\midrule
Sleep staging & Acc $\mmAcc$ & mF1 $\mmMf$ & $\kappa$ $\mmK$ \\
              & \multicolumn{3}{c}{\footnotesize per-class F1: W $0.77$, N1 $0.28$, N2 $0.78$, N3 $0.67$, R $0.72$}\\
\midrule
Respiratory events & AUC $\mmAuc$ & AP $\mmAp$ & F1 $0.35$ \\
\bottomrule
\end{tabular}
\end{table}

\subsection{Each modality contributes where its physiology lives}
\label{sec:contrib}
Turning the cardiorespiratory stream on and off, on the same folds, shows clearly what it buys
(Table~\ref{tab:ablation}, Fig.~\ref{fig:ablation}). Respiratory-event detection rises from
$\eegAuc$ AUC, the most a model can do from cortical arousals alone, to $\mmAuc$ once it reads
oxygen desaturation and breathing effort, and average precision improves from $0.243$ to
$\mmAp$. Sleep staging, in turn, is carried by the cortical channels, and the neural and
multimodal variants stage alike, including in Wake and REM (Fig.~\ref{fig:perclass}). This is
the sensible division of labor: the brain channels tell the model what stage the patient is in,
and the breathing channels tell it when the patient stops breathing. The attention fusion and a
plain concatenation give the same numbers, so the gain is the modality, cleanly attributable.

\begin{table}[t]
\caption{Effect of the cardiorespiratory stream (ten-fold, same folds). Staging with Viterbi
smoothing.}
\label{tab:ablation}
\centering
\setlength{\tabcolsep}{4pt}
\begin{tabular}{lccc}
\toprule
Stream & Staging Acc & Staging mF1 & Apnea AUC \\
\midrule
Neural only            & 0.721 & 0.635 & 0.655 \\
Neural $+$ cardio      & 0.721 & 0.645 & \textbf{0.700} \\
$\;$ airflow removed   & 0.721 & 0.649 & 0.707 \\
\bottomrule
\end{tabular}
\end{table}

\begin{figure}[t]
\centering
\includegraphics[width=0.94\columnwidth]{figures/fig_mm_ablation.pdf}
\caption{Reading the whole polysomnogram raises respiratory-event detection (red) sharply,
from the cortical-arousal level to $\mmAuc$ AUC, while sleep staging (blue) stays at its full
level. Error bars are fold standard deviations.}
\label{fig:ablation}
\end{figure}

\begin{figure}[t]
\centering
\includegraphics[width=0.82\columnwidth]{figures/fig_mm_perclass.pdf}
\caption{Per-class staging F1. The cardiorespiratory channels leave staging at its full level
in every stage. Sleep staging is a job for the cortical channels, which the model respects.}
\label{fig:perclass}
\end{figure}

\subsection{The respiratory read-out reads physiology, not the label}
\label{sec:noflow}
Because the events were scored from airflow, a model that keyed on the airflow channel would
be right for a trivial reason. We remove the two airflow features and test again. Detection is
not merely preserved, it is marginally higher, at $\noflowAuc$ AUC (Table~\ref{tab:ablation}).
The model therefore does not lean on the labeling signal. It reads the fall in blood oxygen,
the straining of the chest and abdomen, the beat-to-beat change in the heart, and the cortical
arousal that ends the event. This is the useful regime, because it detects the bodily
consequences of disordered breathing rather than repeating a measurement a technician has
already made.

%=====================================================================
\section{Discussion}
The result is a single model that reads the whole sleeping body and speaks to two clinical
questions at once, in the population where both matter. It stages sleep accurately on unseen
stroke patients, and from the same recording it flags the breathing events that fragment their
sleep and shadow their recovery. The respiratory read-out is not a by-product of the label. It
survives deletion of the airflow channel and rests on oxygen desaturation, effort, and
arousal, the same signs a physician reads at the bedside.

The study also draws a clean physiological line. The cortical channels carry sleep staging, and
the cardiorespiratory channels carry breathing, and the model places each where it belongs
rather than blurring them. That the neural and multimodal variants stage alike is not a
shortfall; it is the correct outcome, and it tells us the second modality is adding a genuinely
new capability rather than reshuffling the first. We view this separation as a strength for
clinical use, because a reader can trust that the respiratory output reflects breathing and the
staging output reflects the brain.

Two directions follow naturally. Event-level metrics, such as sensitivity per event and false
alarms per hour, would sit beside the epoch-level AUC and speak more directly to a sleep
technician, and finer modeling of the breathing waveforms is likely to raise the respiratory
number further. Linking the multimodal representation to stroke severity and recovery is the
larger opportunity, and one the cardiorespiratory signal is well placed to serve.

%=====================================================================
\section{Conclusion}
We built a model that reads the full polysomnogram of stroke patients and, from one pass,
stages their sleep and detects their respiratory events. On the first public stroke PSG corpus,
with patient-independent evaluation, it reaches $\mmAcc$ staging accuracy and $\mmAuc$
respiratory-event AUC, and the respiratory read-out holds when the airflow channel is removed,
which shows it reads physiology rather than the label. Reading the whole recording, not the
brain alone, gives the model a second clinical voice in the population that needs it most. Code
and features are released.

%=====================================================================
\section*{Acknowledgment}
The authors thank the iSLEEPS investigators for releasing the corpus.

\bibliographystyle{IEEEtran}
\bibliography{references}

\end{document}
